\section{Các số đặc trưng đo xu thế trung tâm cho mẫu số liệu không ghép nhóm}

\subsection{Đề 1}
\Opensolutionfile{ans}[ans/ans-0-B13]
\caulc
\begin{ex}%[0D6N3-5]
Các giá trị xuất hiện nhiều nhất trong mẫu số liệu được gọi là
\choice
{\True Mốt}
{Số trung bình}
{Số trung vị}
{Độ lệch chuẩn}
\loigiai{
Các giá trị xuất hiện nhiều nhất trong mẫu số liệu được gọi là Mốt.
}
\end{ex}

\begin{ex}%[0D6N3-2]
Chiều dài (đơn vị mét) của 5 con cá voi xanh trưởng thành được cho như sau:
$$
\begin{array}{lllll}
25 & 22 & 31 & 33 & 27.
\end{array}
$$
Tìm chiều dài trung bình của $5$ con cá voi xanh trên.
\choice
{$27$ m}
{\True $27{,}6$ m}
{$23$ m}
{$28{,}2$ m}
\loigiai{
Số trung bình của mẫu số liệu $\dfrac{25+22+31+33+27}{5}\approx 27{,}6$ m.}
\end{ex}

\begin{ex}%[0D6H3-2]
Lớp $10$A của một trường trung học phổ thông có điểm thi môn Toán trong học kỳ $1$ vừa qua được thống kê trong bảng sau
\begin{center}
\begin{tabular}{|l|l|l|l|l|l|l|}
\hline Điểm thi & 5 & 6 & 7 & 8 & 9 & 10 \\
\hline Tần số & 2 & 7 & 9 & 12 & 8 & 4 \\
\hline
\end{tabular}
\end{center}
Tính điểm trung bình cộng môn Toán của lớp $10$A (làm tròn đến hàng phần mười).
\choice
{$7{,}6$}
{$7{,}9$}
{\True $7{,}7$}
{$8{,}0$}
\loigiai{
Ta có $\overline{x}=\dfrac{5\times 2+6\times 7+7\times 9+8\times 12+9\times 8+10\times 4}{42}\approx 7{,}7$.}
\end{ex}

\begin{ex}%[0D6H3-2]
Để đủ điều kiện được cấp chứng chỉ A của trung tâm tin học, học viên phải trải qua $6$ lần thi trắc nghiệm, thang điểm mỗi lần là $100$ điểm, và phải đạt trung bình $70$ điểm trở lên. Qua $5$ lần thi bạn Vân đạt trung bình $67{,}5$ điểm. Hỏi trong lần kiểm tra cuối cùng Vân phải đạt ít nhất bao nhiêu điểm để được cấp chứng chỉ?
\choice
{$70$}
{$80$}
{$90$}
{\True $82{,}5$}
\loigiai{
Gọi $x$ là điểm số trong lần kiểm tra cuối mà Vân cần đạt để được cấp chứng chỉ.\\
Ta có số điểm qua $5$ lần thi của Vân là $67{,}5\times 5=337{,}5$.\\
Theo bài ra ta có $\dfrac{x+337{,}5}{6}\geq 70 \Leftrightarrow x \geq 70\times 6-337{,}5 \Leftrightarrow x \geq 82{,}5$.\\
Vậy trong lần kiểm tra cuối cùng, Vân phải đạt ít nhất $82{,}5$ điểm để được cấp chứng chỉ.}
\end{ex}

\begin{ex}%[0D6N3-3]
Cân nặng (kg) của một nhóm học sinh nữ lớp $10$A cho bởi số liệu sau:
$$
\begin{array}{llllllllll}
38 & 39 & 44 & 44 & 46 & 47 & 47 & 48 & 52 & 55.
\end{array}
$$
Tìm số trung vị của mẫu số liệu trên.
\choice
{$44$}
{$46$}
{\True $46{,}5$}
{$47$}
\loigiai{
Dãy trên đã sắp xếp không giảm, và có $10$ số trong mẫu số liệu.\\
Số liệu thứ $5$ là $46$ và số liệu thứ $6$ là $47$.\\
Vậy số trung vị là $M_e=\dfrac{46+47}{2}=46{,}5$.}
\end{ex}

\begin{ex}%[0D6N3-3]
Nhiệt độ $\left({}^{\circ}\mathrm{C}\right)$ trung bình các ngày trong tuần đầu tháng $3$ của Nam Định là
$$
\begin{array}{lllllll}
20 & 19 & 22 & 19 & 24 & 28 & 25
\end{array}
$$
Tìm số trung vị của mẫu số liệu trên.
\choice
{\True $22$}
{$23$}
{$24$}
{$25$}
\loigiai{
Sắp xếp theo thứ tự không giảm: 
$\begin{array}{ccccccc}19 & 19 & 20 & 22 & 24 & 25 & 28.\end{array}$
Mẫu số liệu trên có $7$ số.\\
Số trung vị là số đứng ở vị trí thứ $\dfrac{7+1}{2}=4$.\\
Vậy số trung vị là $M_e=22$}
\end{ex}

\begin{ex}%[0D6H3-3]
Số liệu sau đây cho ta lãi hàng tháng của một cửa hàng trong năm $2023$. Đơn vị là triệu đồng
\begin{center}
\begin{tabular}{|c|c|c|c|c|c|c|c|c|c|c|c|c|}
\hline Tháng & 1 & 2 & 3 & 4 & 5 & 6 & 7 & 8 & 9 & 10 & 11 & 12 \\
\hline Lãi & 10 & 12 & 18 & 15 & 22 & 20 & 15 & 13 & 18 & 17 & 20 & 22 \\
\hline
\end{tabular}
\end{center}
Tìm số trung vị của mẫu số liệu trên.
\choice
{$15$}
{\True $17{,}5$}
{$20$}
{$21$}
\loigiai{
Sắp xếp tiền lãi theo thứ tự không giảm
\begin{center}
\begin{tabular}{|c|c|c|c|c|c|c|c|c|c|c|c|}
\hline 10 & 12 & 13 & 15 & 15 & 17 & 18 & 18 & 20 & 20 & 22 & 22 \\
\hline
\end{tabular}
\end{center}
Mẫu số liệu trên có $12$ số. Số thứ $6$ và số thứ $7$ lần lượt là $17$ và $18$.\\
Vậy số trung vị là $M_e=\dfrac{17+18}{2}=17{,}5$.}
\end{ex}

\begin{ex}%[0D6H3-4]
Tìm tứ phân vị thứ nhất của mẫu số liệu:
\begin{center}
\begin{tabular}{ccccccccccc}
27 & 15 & 18 & 30 & 19 & 40 & 100 & 9 & 46 & 10 & 200.
\end{tabular}
\end{center}
\choice
{$18$}
{\True $15$}
{$40$}
{$46$}
\loigiai{
Sắp xếp số liệu theo thứ tự không giảm: \begin{center}
\begin{tabular}{ccccccccccc}
9 & 10 & 15 & 18 & 19 & 27 & 30 & 40 & 46 & 100 & 200.
\end{tabular}
\end{center}
Tứ phân vị thứ nhì là trung vị của dãy số liệu là $Q_2=27$.\\
Tứ phân vị thứ nhất là trung vị của dãy số liệu: \begin{center}
\begin{tabular}{ccccc}
9 & 10 & 15 & 18 & 19.
\end{tabular}
\end{center}
Khi đó tứ phân vị thứ nhất là $Q_1=15$.}
\end{ex}

\begin{ex}%[0D6H3-4]
Tìm tứ phân vị thứ ba của mẫu số liệu: \begin{center}
\begin{tabular}{ccccccccccc}
27 & 15 & 18 & 30 & 19 & 40 & 100 & 9 & 46 & 10 & 200.
\end{tabular}
\end{center}
\choice
{$18$}
{$15$}
{$40$}
{\True $46$}
\loigiai{
Sắp xếp số liệu theo thứ tự không giảm: \begin{center}
\begin{tabular}{ccccccccccc}
9 & 10 & 15 & 18 & 19 & 27 & 30 & 40 & 46 & 100 & 200.
\end{tabular}
\end{center}
Tứ phân vị thứ nhì là trung vị của dãy số liệu là $Q_2=27$.\\
Tứ phân vị thứ ba là trung vị của dãy số liệu: \begin{center}
\begin{tabular}{ccccc}
30 & 40 & 46 & 100 & 200.
\end{tabular}
\end{center} 
Khi đó tứ phân vị thứ ba là $Q_3=46$.}
\end{ex}

\begin{ex}%[0D6N3-4]
Số học sinh của các lớp khối $10$ trong một trường trung học được thống kê như sau:\begin{center}
\begin{tabular}{cccccccc}
35 & 40 & 37 & 46 & 38 & 42 & 43 & 36
\end{tabular}
\end{center}  Tứ phân vị thứ hai của mẫu số liệu được thống kê ở trên là
\choice
{$40$}
{\True $39$}
{$36$}
{$42$}
\loigiai{
Sắp xếp số liệu theo thứ tự không giảm: \begin{center}
\begin{tabular}{cccccccc}
35 & 36 & 37 & 38 & 40 & 42 & 43 & 46.
\end{tabular}
\end{center} 
Dãy số liệu trên có $2$ số chính giữa nên tứ phân vị thứ hai là $Q_2=\dfrac{38+40}{2}=39$.}
\end{ex}

\begin{ex}%[0D6N3-5]
Số áo bán được trong một quý ở cửa hàng bán áo sơ mi nữ được thống kê như sau:
\begin{center}
\begin{tabular}{|c|l|l|l|l|l|l|l|}
\hline Cỡ áo & 36 & 37 & 38 & 39 & 40 & 41 & 42 \\
\hline 
Tần số 
(Số áo bán được)
& 13 & 45 & 126 & 125 & 110 & 40 & 12 \\
\hline
\end{tabular}
\end{center}
Giá trị mốt của bảng phân bố tần số trên bằng
\choice
{\True $38$}
{$126$}
{$42$}
{$12$}
\loigiai{
Vì giá trị $x_3=38$ có tần số $n_3=126$ lớn nhất nên mốt của bảng số liệu là $M_O=38$.}
\end{ex}

\begin{ex}%[0D6N3-5]
Điều tra tiền lương một tháng của $100$ người lao động trên địa bàn một xã ta có bảng phân bố tần số sau:
\begin{center}
\begin{tabular}{|c|c|c|c|c|c|c|}
\hline 
\begin{tabular}{c}
Tiền lương \\
(VND)
\end{tabular}
& 6 000 000 & 7 000 000 & 8 500 000 & 9 500 000 & 10 000 000 & 10 500 000 \\
\hline Tần số & 13 & 37 & 25 & 11 & 6 & 8 \\
\hline
\end{tabular}
\end{center}
Tìm mốt của bảng phân bố tần số trên.
\choice
{$6\ 000\ 000$}
{\True $7\ 000\ 000$}
{$8\ 500\ 000$}
{$9\ 500\ 000$}
\loigiai{
Ta có giá trị $7\ 000\ 000$ có tần số lớn nhất nên là mốt của bảng phân bố tần số trên.}
\end{ex}
\Closesolutionfile{ans}
\indapan{6}{ans/ans-0-B13}
\cauds
\Opensolutionfile{ans}[ans/ans-0-B13-DS]
\begin{ex}%[0D6H3-5]
Số điểm mà năm vận động viên bóng rổ ghi được trong một trận đấu:
\begin{center}
\begin{tabular}{ccccc}
9 & 8 & 15 & 8 & 20.
\end{tabular}
\end{center}
Các mệnh đề sau đúng hay sai?
\choiceTF
	{\True Số trung bình cộng của mẫu số liệu trên là $12$}
	{Mốt của mẫu số liệu trên là $15$}
	{\True Trung vị là $9$}
	{Tứ phân vị thứ nhất $Q_1=17{,}5$}
\loigiai{
\begin{itemchoice}
\itemch Đúng.  Số trung bình là $\dfrac{8\times 2+9+15+20}{5}=12$.
\itemch Sai. Mốt của mẫu số liệu trên là $8$ vì $8$ có tần số lớn nhất.
\itemch Đúng. Sắp xếp số liệu theo thứ tự không giảm: $8 \quad 8 \quad 9 \quad 15 \quad 20$.\\
Trung vị là $M_e=9$.
\itemch Sai. Tứ phân vị thứ nhất là $Q_1=\dfrac{8+8}{2}=8$.
\end{itemchoice}
}
\end{ex}
\begin{ex}%[0D6H3-5]
Cho bảng số liệu điểm kiểm tra môn Toán của $20$ học sinh.
\begin{center}
\begin{tabular}{|c|c|c|c|c|c|c|c|c|}
\hline Điểm & 4 & 5 & 6 & 7 & 8 & 9 & 10 & Cộng \\
\hline Số học sinh & 1 & 2 & 3 & 4 & 5 & 4 & 1 & 20 \\
\hline
\end{tabular}
\end{center}
Các mệnh đề sau đúng hay sai?
\choiceTF
{Số trung bình $\overline{x}=7{,}5$}
{\True Tứ phân vị thứ hai là $Q_2=7{,}5$}
{\True Tứ phân vị thứ nhất là $Q_1=6$}
{Mốt là $M_O=10$}
\loigiai{
\begin{itemchoice}
\itemch Sai. Số trung bình là $\overline{x}=\dfrac{4+2\times 5+6\times 3+7\times 4+8\times 5+9\times 4+10}{20}=7{,}3$.
\itemch Đúng. Vì cỡ mẫu là $n=20$, là số chẵn, nên giá trị của tứ phân vị thứ hai là $Q_2=\dfrac{7+8}{2}=7{,}5$.
\itemch Đúng. Tứ phân vị thứ nhất là $Q_1=\dfrac{6+6}{2}=6$.
\itemch Sai. Vì $8$ có tần số lớn nhất nên $M_O=8$.
\end{itemchoice}
}
\end{ex}

\begin{ex}%[0D6H3-5]
Thống kê số bao xi măng được bán ra tại một cửa hàng vật liệu xây dựng trong 6 tháng cho kết quả như sau:\begin{center}
\begin{tabular}{|l|l|l|l|l|l|}
\hline 60 & 61 & 65 & 63 & 61 & 71 \\
\hline
\end{tabular}
\end{center}
Các mệnh đề sau đúng hay sai?
\choiceTF
{Mẫu số liệu trên có $n=12$}
{\True Mốt của mẫu số liệu là $61$}
{\True Sai khác giữa số trung bình và số trung vị là $1{,}5$}
{Khoảng cách từ $Q_1$ đến $Q_2$ là $3$}
\loigiai{
\begin{itemchoice}
\itemch Sai. Vì mẫu số liệu trên có $n=6$.
\itemch Đúng. Mốt của mẫu số liệu là $61$, vì số bao xi măng bán ra là $61$ xuất hiện $2$ lần với tần suất lớn nhất.
\itemch Đúng.\\
Ta có số trung bình $\overline{x}=\dfrac{60+61+65+63+61+71}{6}=63{,}5$.\\
Sắp xếp lại dãy số theo mức độ tăng dần ta có dãy
\begin{center}
\begin{tabular}{|r|r|r|r|r|r|}
\hline 60 & 61 & 61 & 63 & 65 & 71 \\
\hline
\end{tabular}
\end{center}
Số trung vị của dãy số liệu là $Q_2=\dfrac{61+63}{2}=62$.\\
Sai khác giữa số trung bình và số trung vị là $1{,}5$.
\itemch Sai. Ta có số trung vị $Q_2=62$.\\
Số trung vị của nửa bên trái $Q_2$ là $Q_1=61$.\\
Khoảng cách từ $Q_1$ đến $Q_2$ là $1$.
\end{itemchoice}
}
\end{ex}
\begin{ex}%[0D6H3-5]
Điểm điều tra về chất lượng sản phẩm mới (thang điểm $100$) như sau:
\begin{center}
\begin{tabular}{lllllllllllll}72 & 68 & 39 & 41 & 54 & 61 & 72 & 75 & 80 & 61 & 50 & 65 & 78\end{tabular}
\end{center}
Các mệnh đề sau đúng hay sai?
\choiceTF 
{$n=12$}
{\True Mốt của mẫu số liệu là $61$}
{Tứ phân vị thứ hai là số thứ $6$}
{Tổng số trung vị, tứ phân vị thứ nhất và tứ phân vị thứ ba là $190$}
\loigiai
{\begin{itemchoice}
\itemch Sai. Mẫu số liệu có $13$ số liệu nên $n=13$.
\itemch Đúng. Vì điểm $61$ xuất hiện $2$  lần với tần suất lớn nhất nên mốt của mẫu số liệu là $61$.
\itemch Sai. Vì $n=13$ là số lẻ nên số trung vị là số đứng ở vị trí thứ $\dfrac{13+1}{2}=7$.
\itemch Sai. Sắp xếp mẫu số liệu theo thứ tự tăng dần ta có
\begin{center}
\begin{tabular}{lllllllllllll}
39 & 41 & 50 & 54 & 61 & 61 & 65 & 68 & 72 & 72 & 75 & 78 & 80.
\end{tabular}
\end{center}
Số trung vị là số đứng thứ $7$ nên $Q_2=65$.\\
Xét nửa mẫu số liệu bên trái: $\begin{array}{llllll}39 & 41 & 50 & 54 & 61 & 61\end{array}$.\\
Tứ phân vị thứ nhất chính là trung vị của mẫu số liệu này.\\ Như thế $Q_1=\dfrac{50+54}{2}=52$.\\
 Xét nửa mẫu số liệu bên phải: $\begin{array}{llllll}68 & 72 & 72 & 75 & 78 & 80\end{array}$.\\
Tứ phân vị thứ ba chính là trung vị nửa mẫu số liệu này. Do đó $Q_3=\dfrac{72+75}{2}=73{,}5$.\\
Tổng số trung vị, tứ phân vị thứ nhất và tứ phân vị thứ ba là $65+52+73{,}5=190{,}5$.
\end{itemchoice}
}
\end{ex}

\Closesolutionfile{ans}
\indapan{3}{ans/ans-0-B13-DS}

\Opensolutionfile{ans}[ans/ans-0-B13-KQ]
\caukq
\begin{ex}%[0D6H3-5]
Thời gian (đơn vị giờ) dành cho hoạt động thể thao trong tuần của một số học sinh được thống kê như sau:
\begin{center}
\begin{tabular}{|l|l|l|l|l|l|l|l|l|l|}
\hline 1 & 0 & 3 & 2 & 1 & 1 & 4 & 1,5 & 2 & 4 \\
\hline
\end{tabular}
\end{center}
Mốt của mẫu số liệu trên bằng bao nhiêu?
	\shortans{$1$}
\loigiai{
Vì số học sinh dành $1$ giờ mỗi tuần cho hoạt động thể thao là nhiều nhất ($3$ học sinh) nên mốt là $1$.
}
\end{ex}

\begin{ex}%[0D6N3-5]
Có $100$ học sinh tham gia kì thi khảo sát chất lượng môn Toán. Điểm khảo sát được tính theo thang điểm $10$ và thống kê như sau:
\begin{center}
\begin{tabular}{|l|l|l|l|l|l|l|l|l|l|l|}
\hline Điểm & 1 & 2 & 3 & 4 & 5 & 6 & 7 & 8 & 9 & 10 \\
\hline Số & 11 & 8 & 41 & 56 & 29 & 34 & 61 & 50 & 91 & 19 \\
\hline
\end{tabular}
\end{center}
Mốt của mẫu số liệu trên bằng bao nhiêu?
	\shortans{$9$}
\loigiai{
Số học sinh đạt điểm $9$ là nhiều nhất ($91$ học sinh) nên mốt của mẫu số liệu là $9$.
}
\end{ex}

\begin{ex}%[0D6H3-3]
Cho mẫu số liệu: $2\quad 3\quad 7\quad 4\quad 5\quad 3\quad 4\quad 4$. Số trung vị của mẫu số trên bằng bao nhiêu?
	\shortans{$4$}
\loigiai{
Sắp xếp mẫu: $2\quad 3\quad 3\quad 4\quad 4\quad 4\quad 5\quad 7$.\\
Kích thước mẫu là $8$.
Khi đó trung vị của mẫu là trung bình cộng của giá trị thứ $4$ và $5$ là $M_e=\dfrac{1}{2}\cdot(4+4)=4$.
}
\end{ex}

\begin{ex}%[0D6H3-3]
Một cửa hàng giày thể thao đã thống kê cỡ giày của $20$ khách hàng nữ được chọn ngẫu nhiên cho kết quả như sau:
$$
35\quad 37\quad 39\quad 41\quad 38\quad 40\quad 40\quad 37\quad 39\quad 38\quad 38\quad 36\quad 37\quad 42\quad 38\quad 35\quad 38\quad 36\quad 38\quad 35.
$$
Số trung vị của mẫu số trên bằng bao nhiêu?
	\shortans{$38$}
\loigiai{
Sắp xếp các giá trị theo thứ tự không giảm:
$$
35\quad 35\quad 35\quad 36\quad 36\quad 37\quad 37\quad 37\quad 38\quad 38\quad 38\quad 38\quad 38\quad 38\quad 39\quad 39\quad 40\quad 40\quad 41\quad 42 .
$$
Vì $n=20$ là số chẵn nên trung vị là trung bình cộng của hai giá trị chính giữa
$$
M_e=\dfrac{38+38}{2}=38.$$}
\end{ex}

\begin{ex}%[0D6H3-3]
Bảng sau cho biết thời gian chạy cự li $100$ m của các bạn trong lớp (đơn vị giây)
\begin{center}
\begin{tabular}{|l|l|l|l|l|l|}
\hline Thời gian & 12 & 13 & 14 & 15 & 16 \\
\hline Số bạn & 4 & 7 & 3 & 18 & 8 \\
\hline
\end{tabular}
\end{center}
Số tứ phân vị $Q_3$ của mẫu số liệu trên bằng bao nhiêu?
	\shortans{$15$}
\loigiai{
Số bạn học sinh trong lớp là $n=4+7+3+18+8=40$ (bạn).\\
Tứ vị phân thứ ba là $Q_3=\dfrac{15+15}{2}=15$.}
\end{ex}


\begin{ex}%[0D6H2-3]
Một lớp gồm $40$ học sinh chia đều thành $4$ tổ. Nhân dịp Tết Nguyên Đán, các em cùng đóng góp để giúp đỡ các gia đình khó khăn cùng đón cái Tết đầm ấm. Biết rằng mỗi em đóng góp $10$ đến $15$ ngàn đồng. Cuối đợt đóng góp, lớp trưởng thống kê lại như bảng sau 
\begin{center}
\begin{tabular}{|c|c|c|c|c|}
\hline Tổ& 1 & 2 & 3 & 4 \\
\hline Số tiền (đồng) & 90 000 & 102 000 & 200 000 & 140 000 \\
\hline
\end{tabular}
\end{center}
Hỏi lớp trưởng thống kê sai mấy tổ?
\shortans{$2$}
\loigiai{
Số học sinh mỗi tổ là $\dfrac{40}{4}=10$ học sinh.\\
Số tiển tối thiểu mỗi tổ quyên góp là $10\times 10\ 000=100\ 000$ đồng.\\
Số tiền tối đa mỗi tổ quyên góp là $10\times 15\ 000=150\ 000$ đồng.\\
Vậy lớp trưởng thống kê sai $2$ tổ (tổ $1$ và tổ $3$).}
\end{ex}
\Closesolutionfile{ans}
\indapan{6}{ans/ans-0-B13-KQ}
\subsection{Đề 2}
\Opensolutionfile{ans}[ans/ans-0-B15]
\caulc
\begin{ex}%[0D6N4-1]
	Để đánh giá mức độ phân tán của các số liệu thống kê so với số trung bình, ta dùng đại lượng nào sau đây?
	\choice
	{Số trung bình}
	{Số trung vị}
	{Mốt}
	{\True Phương sai}
	\loigiai{
		Dựa vào ý nghĩa của phương sai và độ lệch chuẩn để đo mức độ phân tán của các số liệu trong mẫu quanh số trung bình.
	}
\end{ex}	

\begin{ex}%[0D6N4-1]
	Chọn câu đúng trong các câu trả lời sau đây: Phương sai bằng
	\choice
	{một nửa của độ lệch chuẩn}
	{căn bậc hai của độ lệch chuẩn}
	{hai lần của độ lệch chuẩn}
	{\True bình phương của độ lệch chuẩn}
	\loigiai{
		Ta có
		\begin{itemize}
			\item Phương sai là $s_x^2$.
			\item Độ lệch chuẩn là $s_x=\sqrt{s_x^2}$.
		\end{itemize}
		Suy ra phương sai bằng bình phương của độ lệch chuẩn.
	}
\end{ex}	

\begin{ex}%[0D6N4-4]
	Cho phương sai của các số liệu bằng $4$. Tìm độ lệch chuẩn.
	\choice
	{$4$}
	{\True $2$}
	{$16$}
	{$8$}
	\loigiai{
		Ta có độ lệch chuẩn là căn bậc hai của phương sai. Suy ra \[ s_x=\sqrt{s_x^2}=\sqrt{4}=2.\]
	}
\end{ex}	

\begin{ex}%[0D6N4-4]
	Độ lệch chuẩn là
	\choice
	{\True căn bậc hai của phương sai}
	{bình phương của phương sai}
	{một nửa của phương sai}
	{không phải các công thức trên}
	\loigiai{
		Ta có độ lệch chuẩn là $s_x=\sqrt{s_x^2}$ nên độ lệch chuẩn là căn bậc hai của phương sai.
	}
\end{ex}	

\begin{ex}%[0D6N4-1]
	Tìm phát biểu đúng về phương sai của một mẫu số liệu.
	\choice
	{Phương sai được sử dụng làm đại diện cho các số liệu của mẫu}
	{\True Phương sai được sử dụng để đánh giá mức độ phân tán của các số liệu thống kê (so với số trung bình)}
	{Phương sai được tính bằng tổng số phần tử của một mẫu số liệu}
	{Phương sai là số liệu xuất hiện nhiều nhất (số liệu có tần số lớn nhất) trong bảng các số
		liệu thống kê}
	\loigiai{
		Ta có ý nghĩa của phương sai là \lq\lq Phương sai được sử dụng để đánh giá mức độ phân tán của các số liệu thống kê (so với số trung bình).\rq\rq
	}
\end{ex}	

\begin{ex}%[0D6N4-4]
	Theo kết quả thống kê điểm thi giữa kỳ 2 môn toán khối 11 của một trường THPT, người ta tính được phương sai của bảng thống kê đó là $s_x^2=0{,}573$. Độ lệch chuẩn của bảng thống kê đó xấp xỉ bằng
	\choice
	{$0{,}812$}
	{\True $0{,}757$}
	{$0{,}936$}
	{$0{,}657$}
	\loigiai{
		Ta có \[ s_x=\sqrt{s_x^2}=\sqrt{0{,}573}\approx 0{,}757.\]
	}
\end{ex}	

\begin{ex}%[0D6N4-4]
	Phương sai của dãy số $2$; $3$; $4$; $5$; $6$ là
	\choice
	{$s_x^2=4$}
	{$s_x^2=\sqrt{2}$}
	{\True $s_x^2=2$}
	{$s_x^2=-2$}
	\loigiai{
		Ta có $\overline{x}=\dfrac{2+3+4+5+6}{5}=4$ nên suy ra \[s_x^2=\dfrac{1}{5}\left[ (2-4)^2+(3-4)^2+(5-4)^2+(6-4)^2\right]=2.\]
	}
\end{ex}

\begin{ex}%[0D6N4-2]
	Khoảng tứ phân vị của dãy số $2$; $3$; $4$; $5$; $6$ là
	\choice
	{\True $\Delta_Q=3$}
	{$\Delta_Q=\sqrt{2}$}
	{$\Delta_Q=2$}
	{$\Delta_Q=-2$}
	\loigiai{
		Ta có \[\Delta_Q=Q_3-Q_1=\dfrac{11}{2}-\dfrac{5}{2}=3.\]
	}
\end{ex}	

\begin{ex}%[0D6N4-2]
	Cho mẫu số liệu thống kê $\{ 1;2;3;4;5;6;7;8;9 \}$. Tìm khoảng tứ phân vị của mẫu số liệu trên.
	\choice
	{$2$}
	{\True $5$}
	{$3$}
	{$4$}
	\loigiai{
		Ta có $Q_1=2{,}5$; $Q_3=7{,}5$ nên suy ra \[ \Delta_Q=Q_3-Q_1=5.\]
	}
\end{ex}	

\begin{ex}%[0D6N4-2]
	Cho dãy số liệu thống kê $1$; $2$; $3$; $4$; $5$; $6$; $7$. Khoảng biến thiên là 
	\choice
	{$1$}
	{$2$}
	{$3$}
	{\True $6$}
	\loigiai{
		Ta có $7-1=6$.
	}
\end{ex}	

\begin{ex}%[0D6H4-4]
	Điểm thi của lớp 10C của một trường Trung học Phổ Thông được trình bày ở bảng phân bố tần số sau
	\begin{center}
		\begin{tabular}{|c|c|c|c|c|c|c|c|}
			\hline Điểm thi	& $5$ & $6$ & $7$ & $8$ & $9$ & $10$ & \\
			\hline Tần số	& $7$ & $5$ & $10$ & $12$ & $4$ & $2$ & $n=40$\\
			\hline 
		\end{tabular}
	\end{center}
	Phương sai của bảng phân bố tần số đã cho là 
	\choice
	{$0{,}94$}
	{$3{,}94$}
	{$2{,}94$}
	{\True $1{,}94$}
	\loigiai{
		Trong dãy số liệu về điểm thi của lớp 10C ta có
		\[ \overline{x} = \dfrac{1}{n}\left( n_1x_1+n_2x_2+ \cdots + n_6x_6\right) =\dfrac{1}{40}\left(7\cdot 5+5\cdot 6+10\cdot 7+12\cdot 8+4\cdot 9+2\cdot 10\right)=7{,}175.\]
		Suy ra phương sai là
		\begin{align*}
			s^2 &= \dfrac{1}{n}\left[ n_1\left(x_1-\overline{x}\right)^2 + n_2\left(x_2-\overline{x}\right)^2 + \cdots + n_6\left(x_6-\overline{x}\right)^2\right] 
			\\ &=  \dfrac{1}{40}\left[ 7\left(5-7{,}175\right)^2 + 5\left(6-7{,}175\right)^2 + 10\left(7-7{,}175\right)^2 + 12\left(8-7{,}175\right)^2 \right.
			\\ & \quad \left. + 4\left(9-7{,}175\right)^2 + 2\left(10-7{,}175\right)^2 \right] 
			\\ & \approx 1{,}94.
		\end{align*}
	}
\end{ex}	

\begin{ex}%[0D6H4-4]
	Theo dõi thời gian làm một bài toán (tính bằng phút) của $40$ học sinh, giáo viên lập được bảng sau
	\begin{center}
		\begin{tabular}{|c|c|c|c|c|c|c|c|c|c|c|}
			\hline Thời gian $(x)$	& $4$ & $5$ & $6$ & $7$ & $8$ & $9$ & $10$ & $11$ & $12$ & \\
			\hline Tần số $(n)$		& $6$ & $3$ & $4$ & $2$ & $7$ & $5$ & $5$ & $7$ & $1$ & $N=40$\\
			\hline 
		\end{tabular}
	\end{center}
	Phương sai của mẫu số liệu trên gần với số nào nhất? 
	\choice
	{\True $6$}
	{$12$}
	{$40$}
	{$9$}
	\loigiai{
		Ta có giá trị trung bình của mẫu số liệu là
		\[ \overline{x} = \dfrac{1}{N}\left( n_1x_1+n_2x_2+ \cdots + n_kx_k\right) =\dfrac{317}{40}.\]
		Suy ra phương sai của mẫu số liệu là
		\[ s^2 = \dfrac{1}{N}\left[ n_1\left(x_1-\overline{x}\right)^2 + n_2\left(x_2-\overline{x}\right)^2 + \cdots + n_k\left(x_k-\overline{x}\right)^2\right]  \approx 6.\]
	}
\end{ex}	

\Closesolutionfile{ans}
\indapan{6}{ans/ans-0-B15}

\cauds
\Opensolutionfile{ans}[ans/ans-0-B15-DS]
\begin{ex}%[0D6H4-2]
	Khoảng biến thiên tổng số giờ nắng trong năm của một tỉnh thành được thống kê từ năm 2006 đến 2019 được cho như sau
	\begin{center}
		\begin{tabular}{|c|c|c|c|c|c|c|}
			\hline $1 \, 884$ & $1\, 600$ & $1\, 645$ & $2\, 049{,}9$ & $1\, 913{,}8$ & $1\, 664{,}1$ & $1\, 846{,}5$  \\
			\hline $1\, 964{,}8$ & $1\, 951$ & $2\, 023{,}6$ & $1\, 996{,}2$ & $1\, 699{,}1$ & $1\, 845$ & $2\, 190{,}4$ \\
			\hline 
		\end{tabular}
	\end{center}
	Các mệnh đề sau đúng hay sai?
	\choiceTF
	{Số giờ nắng trung bình trong năm là $1\, 826{,}67$ giờ}
	{\True Số giờ nắng nhỏ nhất $1\, 600$ giờ}
	{\True Số giờ nắng lớn nhất là $2\, 190{,}4$ giờ}
	{Khoảng biến thiên là $520{,}4$}
	\loigiai{
		\begin{itemchoice}
			\itemch Sai. Số giờ nắng trung bình là \[\overline{x}=\dfrac{x_1+x_2+\cdots+x_{14}}{14}=1\, 876{,}67 \text{ giờ}.\]
			\itemch Đúng. Từ bảng số liệu ta có số giờ nắng nhỏ nhất là $1\, 600$ giờ. 
			\itemch Đúng. Từ bảng số liệu ta có số giờ nắng lớn nhất $2\, 190{,}4$ giờ.
			\itemch Sai. Khoảng biến thiên là \[2\, 190{,}4-1\, 600=590{,}4.\]
		\end{itemchoice}
	}
\end{ex}

\begin{ex}%[0D6H4-4]
	Số liệu thống kê tỉ lệ $(\%)$ tốt nghiệp THPT của một địa phương từ năm học 2001 - 2002 đến năm học 2016 - 2017 được cho như sau
	\begin{center}
		\begin{tabular}{|c|c|c|c|c|c|c|c|}
			\hline $98{,}82$ & $97{,}46$ & $99{,}19$ & $98{,}90$ & $98{,}65$ & $79{,}51$ & $85{,}06$ & $86{,}18$  \\
			\hline $98{,}68$ & $99{,}23$ & $99{,}93$ & $99{,}34$ & $99{,}74$ & $93{,}08$ & $97{,}34$ & $97{,}82$ \\
			\hline 
		\end{tabular}
	\end{center}
	Các mệnh đề sau đúng hay sai?
	\choiceTF
	{\True Tỉ lệ tốt nghiệm trung bình $95{,}56 \%$}
	{$99{,}19$ là tỉ lệ $(\%)$ tốt nghiệp THPT cao nhất}
	{Phương sai $s^2\approx 36{,}03$}
	{\True Độ lệch chuẩn $s \approx 6{,}09$}
	\loigiai{
		\begin{itemchoice}
			\itemch Đúng. Tỉ lệ tốt nghiệp trung bình là \[\overline{x}=\dfrac{x_1+x_2+\cdots+x_{16}}{16}=95{,}56 \%.\]
			\itemch Sai. Tỉ lệ tốt nghiệp THPT cao nhất là $99{,}74 \%$.
			\itemch Sai. Phương sai là \[ s^2=\dfrac{1}{16}\left[ \left(x_1-\overline{x}\right)^2+ \left(x_2-\overline{x}\right)^2+\cdots +  \left(x_{16}-\overline{x}\right)^2\right]\approx 37{,}03.\]
			\itemch Đúng. Độ lệch chuẩn là \[s=\sqrt{s^2}\approx 6{,}09.\]
		\end{itemchoice}
	}
\end{ex}

\begin{ex}%[0D6H4-2]
	Cho mẫu số liệu $15$; $20$; $1$; $2$; $4$; $3$; $7$; $5$.
	\\ Các mệnh đề sau đúng hay sai?
	\choiceTF
	{\True Cỡ mẫu là $8$}
	{Khoảng biến thiên của mẫu số liệu là $10$}
	{\True Mẫu số liệu có giá trị trung vị là $4{,}5$}
	{Khoảng tứ phân vị của mẫu là $8$}
	\loigiai{
		Sắp xếp mẫu số liệu theo thứ tự không giảm ta được
		\[1;2;3;4;5;7;15;20.\]
		\begin{itemchoice}
			\itemch Đúng. Từ mẫu số liệu ta có cỡ mẫu là $8$.
			\itemch Sai. Khoảng biến thiên là \[ 20-1=19.\]
			\itemch Đúng. Cỡ mẫu $n=8$ là số chẵn nên trung vị là \[M_e=\dfrac{1}{2}(4+5)=4{,}5.\]
			\itemch Sai. Tứ phân vị thứ nhất và thứ ba của mẫu lần lượt là
			\[Q_1=\dfrac{1}{2}(2+3)=2{,}5 ; Q_3=\dfrac{1}{2}(7+15)=11.\]
			Suy ra khoảng tứ phân vị của mẫu là 
			\[D_Q=Q_3-Q_1=8{,}5.\]
		\end{itemchoice}
	}
\end{ex}

\begin{ex}%[0D6H4-4]
	Điều tra thu nhập của $40$ hộ gia đình ở một bản được số liệu như sau
	\begin{center}
		\begin{tabular}{|c|c|c|c|c|c|c|c|c|c|}
			\hline Thu nhập (triệu đồng/năm) & $4$ & $4{,}5$ & $5$ & $5{,}5$ & $6$ & $6{,}5$ & $7$ & $7{,}5$ & $8$  \\
			\hline Số hộ gia đình & $1$ & $3$ & $4$ & $6$ & $8$ & $7$ & $6$ & $3$ & $2$ \\
			\hline 
		\end{tabular}
	\end{center}
	Các mệnh đề sau đúng hay sai?
	\choiceTF
	{\True Cỡ mẫu là $40$}
	{Mốt của mẫu số liệu là $8$}
	{\True Số trung bình của mẫu số liệu trên là $6{,}1125$}
	{Phương sai của mẫu số liệu trên là $1{,}9561$}
	\loigiai{
		\begin{itemchoice}
			\itemch Đúng. Từ giả thiết ta có cỡ mẫu là $40$.
			\itemch Sai. Có $8$ hộ gia đình thu nhập $6$ triệu đồng/năm nên mốt là $6$.
			\itemch Đúng. Số trung bình của mẫu số liệu trên là \[\overline{x}=\dfrac{4\cdot 1+4{,}5\cdot 3+5\cdot 4+5{,}5\cdot 6+6\cdot 8+6{,}5\cdot 7+7\cdot 6+7{,}5\cdot 3 +8\cdot 2}{40}=6{,}1125.\]
			\itemch Sai. Phương sai của mẫu số liệu trên là 
			\begin{align*}
				s^2 &= \dfrac{1}{40}\left(4^2\cdot 1+4{,}5^2\cdot 3+5^2\cdot 4+5{,}5^2\cdot 6+6^2\cdot 8+6{,}5^2\cdot 7+7^2\cdot 6+7{,}5^2\cdot 3 +8^2\cdot 2\right)-6{,}1125^2
				\\ &\approx 0{,}9561.
			\end{align*}
		\end{itemchoice}
	}
\end{ex}


\Closesolutionfile{ans}
\indapan{3}{ans/ans-0-B15-DS}

\Opensolutionfile{ans}[ans/ans-0-B15-KQ]
\caukq
\begin{ex}%[0D6H4-4]
	Cho mẫu số liệu thống kê $\{1;2;3;4;5;6;7;8;9\}$. Tính độ lệch chuẩn của mẫu số liệu trên (làm tròn đến hàng phần trăm).
	\shortans{$2{,}58$}
	\loigiai{
		Ta có giá trị trung bình của mẫu số liệu là \[\overline{x}=\dfrac{1+2+3+4+5+6+7+8+9}{9}=5.\]
		Suy ra độ lệch chuẩn là \[s=\sqrt{\dfrac{1}{9}\left(1^2+2^2+3^2+4^2+5^2+6^2+7^2+8^2+9^2\right)-5^2}\approx 2{,}58.\]
	}
\end{ex}

\begin{ex}%[0D6H4-4]
	Sản lượng lúa (đơn vị là tạ) của $40$ thửa ruộng thí nghiệm có cùng diện tích được trình bày trong bảng số liệu sau
	\begin{center}
		\begin{tabular}{|c|c|c|c|c|c|c|c|c|c|}
			\hline Sản lượng & $20$ & $21$ & $22$ & $23$ & $24$ &  \\
			\hline Tần số & $5$ & $8$ & $11$ & $10$ & $6$ & $N=40$  \\
			\hline 
		\end{tabular}
	\end{center}
	Tìm phương sai của  mẫu số liệu trên (làm tròn đến hàng phần trăm).
	\shortans{$1{,}54$}
	\loigiai{
		Ta có giá trị trung bình của mẫu số liệu là \[\overline{x}=\dfrac{20\cdot 5+21\cdot 8+22\cdot 11+23\cdot 10+24\cdot 6}{40}=22{,}1.\]
		Suy ra phương sai là \[s^2=\dfrac{1}{40}\left(20^2\cdot 5+21^2\cdot 8+22^2\cdot 11+23^2\cdot 10+24^2\cdot 6\right)-22{,}1^2\approx 1{,}54.\]
	}
\end{ex}

\begin{ex}%[0D6H4-3]
	Cho mẫu số liệu thống kê sau
	\[ -3 \quad 5 \quad 10 \quad 12 \quad 14 \quad 18 \quad 24 \quad 26 \quad 49 \quad 60\]
	Trong các giá trị của mẫu số liệu trên, hãy tìm giá trị ngoại lệ của mẫu số liệu.
	\shortans{$60$}
	\loigiai{
		Từ bảng số liệu ta tìm được số trung vị $Q_2=\dfrac{14+18}{2}=16$, tứ phân vị thứ nhất $Q_1=10$, tứ phân vị thứ ba $Q_3=26$ và khoảng tứ phân vị $\Delta_Q=26-10=16$. Ta có \[\left[ Q_1-1{,}5\Delta_Q;Q_3+1{,}5\Delta_Q\right]=\left[ -14;50\right].\]
		Từ đó ta có $60$ là số liệu ngoại lệ duy nhất của mẫu số liệu.
	}
\end{ex}

\begin{ex}%[0D6H4-4]
	Trong $5$ lần nhảy xa, hai bạn Hùng và Trung có kết quả (đơn vị mét) lần lượt là
	\begin{center}
		\begin{tabular}{|c|c|c|c|c|c|c|c|c|}
			\hline Hùng & $2{,}4$ & $2{,}6$ & $2{,}4$ & $2{,}5$ & $2{,}6$   \\
			\hline Trung & $2{,}4$ & $2{,}5$ & $2{,}5$ & $2{,}5$ & $2{,}6$  \\
			\hline 
		\end{tabular}
	\end{center}
	Hỏi phương sai mẫu số liệu của Hùng gấp của Trung bao nhiêu lần?
	\shortans{$2$}
	\loigiai{
		Phương sai của mẫu số liệu là
		\begin{itemize}
			\item Hùng $s^2=\dfrac{1}{5}\left(2{,}4^2+2{,}6^2+2{,}4^2+2{,}5^2+2{,}6^2\right)-2{,}5^2\approx 0{,}008$.
			\item Trung $s^2=\dfrac{1}{5}\left(2{,}4^2+2{,}5^2+2{,}5^2+2{,}5^2+2{,}6^2\right)-2{,}5^2\approx 0{,}004$.
		\end{itemize}
		Vậy phương sai mẫu số liệu của Hùng gấp của Trung $2$ lần.
	}
\end{ex}

\begin{ex}%[0D6H4-4]
	Bảng số liệu sau thống kê nhiệt độ tại Thành phố Hồ Chí Minh trong một lần đo vào một ngày của năm 2021 như sau
	\begin{center}
		\begin{tabular}{|c|c|c|c|c|c|c|c|c|c|c|c}
			\hline Giờ đo & $1$h & $4$h & $7$h & $10$h & $13$h & $16$h & $19$h & $22$h   \\
			\hline Nhiệt độ ($^\circ \mathrm{C}$) & $27$ & $26$ & $28$ & $32$ & $34$ & $35$ & $30$ & $28$  \\
			\hline 
		\end{tabular}
	\end{center}
	Tìm độ lệch chuẩn của mẫu số liệu đã cho (làm tròn kết quả đến hàng phần trăm).
	\shortans{$3{,}12$}
	\loigiai{
		Ta có số trung bình là \[\overline{x}=\dfrac{27+26+28+32+34+35+30+28}{8}=30^\circ C.\]
		Suy ra độ lệch chuẩn là \[s=\sqrt{\dfrac{1}{8}\left[ (27-30)^2+(26-30)^2+(28-30)^2+\cdots+(28-30)^2\right] }\approx 3{,}12.\]
	}
\end{ex}

\begin{ex}%[0D6H4-4]
	Để biết cây đậu phát triển như thế nào sau khi gieo hạt, bạn Châu gieo $5$ hạt đậu vào $5$ chậu riêng biệt và cung cấp cho chúng lượng nước, ánh sáng như nhau. Sau $2$ tuần, $5$ hạt đậu đã nảy mầm và phát triển thành $5$ cây con. Bạn Châu đo chiều cao từ rễ đến ngọn của mỗi cây (đơn vị mm) và ghi kết quả là mẫu số liệu sau
	\[ 112 \quad 102 \quad 106 \quad 94 \quad 101 \]
	Tìm độ lệch chuẩn của mẫu số liệu trên (kết quả làm tròn đến hàng phần trăm).
	\shortans{$5{,}93$}
	\loigiai{
		Từ bảng số liệu ta có số trung bình là \[ \overline{x}=\dfrac{112+102+106+94+101}{5}=103.\]
		Suy ra ta tính được phương sai và độ lệch chuẩn lần lượt là \[s^2=35{,}2 \text{ và } s\approx 5{,}93.\]
	}
\end{ex}

\Closesolutionfile{ans}
\indapan{6}{ans/ans-0-B15-KQ}
\subsection{Đề 2}
\Opensolutionfile{ans}[ans/ans-0-B15]
\caulc
\begin{ex}%[0D6N4-1]
	Để đánh giá mức độ phân tán của các số liệu thống kê so với số trung bình, ta dùng đại lượng nào sau đây?
	\choice
	{Số trung bình}
	{Số trung vị}
	{Mốt}
	{\True Phương sai}
	\loigiai{
		Dựa vào ý nghĩa của phương sai và độ lệch chuẩn để đo mức độ phân tán của các số liệu trong mẫu quanh số trung bình.
	}
\end{ex}	

\begin{ex}%[0D6N4-1]
	Chọn câu đúng trong các câu trả lời sau đây: Phương sai bằng
	\choice
	{một nửa của độ lệch chuẩn}
	{căn bậc hai của độ lệch chuẩn}
	{hai lần của độ lệch chuẩn}
	{\True bình phương của độ lệch chuẩn}
	\loigiai{
		Ta có
		\begin{itemize}
			\item Phương sai là $s_x^2$.
			\item Độ lệch chuẩn là $s_x=\sqrt{s_x^2}$.
		\end{itemize}
		Suy ra phương sai bằng bình phương của độ lệch chuẩn.
	}
\end{ex}	

\begin{ex}%[0D6N4-4]
	Cho phương sai của các số liệu bằng $4$. Tìm độ lệch chuẩn.
	\choice
	{$4$}
	{\True $2$}
	{$16$}
	{$8$}
	\loigiai{
		Ta có độ lệch chuẩn là căn bậc hai của phương sai. Suy ra \[ s_x=\sqrt{s_x^2}=\sqrt{4}=2.\]
	}
\end{ex}	

\begin{ex}%[0D6N4-4]
	Độ lệch chuẩn là
	\choice
	{\True căn bậc hai của phương sai}
	{bình phương của phương sai}
	{một nửa của phương sai}
	{không phải các công thức trên}
	\loigiai{
		Ta có độ lệch chuẩn là $s_x=\sqrt{s_x^2}$ nên độ lệch chuẩn là căn bậc hai của phương sai.
	}
\end{ex}	

\begin{ex}%[0D6N4-1]
	Tìm phát biểu đúng về phương sai của một mẫu số liệu.
	\choice
	{Phương sai được sử dụng làm đại diện cho các số liệu của mẫu}
	{\True Phương sai được sử dụng để đánh giá mức độ phân tán của các số liệu thống kê (so với số trung bình)}
	{Phương sai được tính bằng tổng số phần tử của một mẫu số liệu}
	{Phương sai là số liệu xuất hiện nhiều nhất (số liệu có tần số lớn nhất) trong bảng các số
		liệu thống kê}
	\loigiai{
		Ta có ý nghĩa của phương sai là \lq\lq Phương sai được sử dụng để đánh giá mức độ phân tán của các số liệu thống kê (so với số trung bình).\rq\rq
	}
\end{ex}	

\begin{ex}%[0D6N4-4]
	Theo kết quả thống kê điểm thi giữa kỳ 2 môn toán khối 11 của một trường THPT, người ta tính được phương sai của bảng thống kê đó là $s_x^2=0{,}573$. Độ lệch chuẩn của bảng thống kê đó xấp xỉ bằng
	\choice
	{$0{,}812$}
	{\True $0{,}757$}
	{$0{,}936$}
	{$0{,}657$}
	\loigiai{
		Ta có \[ s_x=\sqrt{s_x^2}=\sqrt{0{,}573}\approx 0{,}757.\]
	}
\end{ex}	

\begin{ex}%[0D6N4-4]
	Phương sai của dãy số $2$; $3$; $4$; $5$; $6$ là
	\choice
	{$s_x^2=4$}
	{$s_x^2=\sqrt{2}$}
	{\True $s_x^2=2$}
	{$s_x^2=-2$}
	\loigiai{
		Ta có $\overline{x}=\dfrac{2+3+4+5+6}{5}=4$ nên suy ra \[s_x^2=\dfrac{1}{5}\left[ (2-4)^2+(3-4)^2+(5-4)^2+(6-4)^2\right]=2.\]
	}
\end{ex}

\begin{ex}%[0D6N4-2]
	Khoảng tứ phân vị của dãy số $2$; $3$; $4$; $5$; $6$ là
	\choice
	{\True $\Delta_Q=3$}
	{$\Delta_Q=\sqrt{2}$}
	{$\Delta_Q=2$}
	{$\Delta_Q=-2$}
	\loigiai{
		Ta có \[\Delta_Q=Q_3-Q_1=\dfrac{11}{2}-\dfrac{5}{2}=3.\]
	}
\end{ex}	

\begin{ex}%[0D6N4-2]
	Cho mẫu số liệu thống kê $\{ 1;2;3;4;5;6;7;8;9 \}$. Tìm khoảng tứ phân vị của mẫu số liệu trên.
	\choice
	{$2$}
	{\True $5$}
	{$3$}
	{$4$}
	\loigiai{
		Ta có $Q_1=2{,}5$; $Q_3=7{,}5$ nên suy ra \[ \Delta_Q=Q_3-Q_1=5.\]
	}
\end{ex}	

\begin{ex}%[0D6N4-2]
	Cho dãy số liệu thống kê $1$; $2$; $3$; $4$; $5$; $6$; $7$. Khoảng biến thiên là 
	\choice
	{$1$}
	{$2$}
	{$3$}
	{\True $6$}
	\loigiai{
		Ta có $7-1=6$.
	}
\end{ex}	

\begin{ex}%[0D6H4-4]
	Điểm thi của lớp 10C của một trường Trung học Phổ Thông được trình bày ở bảng phân bố tần số sau
	\begin{center}
		\begin{tabular}{|c|c|c|c|c|c|c|c|}
			\hline Điểm thi	& $5$ & $6$ & $7$ & $8$ & $9$ & $10$ & \\
			\hline Tần số	& $7$ & $5$ & $10$ & $12$ & $4$ & $2$ & $n=40$\\
			\hline 
		\end{tabular}
	\end{center}
	Phương sai của bảng phân bố tần số đã cho là 
	\choice
	{$0{,}94$}
	{$3{,}94$}
	{$2{,}94$}
	{\True $1{,}94$}
	\loigiai{
		Trong dãy số liệu về điểm thi của lớp 10C ta có
		\[ \overline{x} = \dfrac{1}{n}\left( n_1x_1+n_2x_2+ \cdots + n_6x_6\right) =\dfrac{1}{40}\left(7\cdot 5+5\cdot 6+10\cdot 7+12\cdot 8+4\cdot 9+2\cdot 10\right)=7{,}175.\]
		Suy ra phương sai là
		\begin{align*}
			s^2 &= \dfrac{1}{n}\left[ n_1\left(x_1-\overline{x}\right)^2 + n_2\left(x_2-\overline{x}\right)^2 + \cdots + n_6\left(x_6-\overline{x}\right)^2\right] 
			\\ &=  \dfrac{1}{40}\left[ 7\left(5-7{,}175\right)^2 + 5\left(6-7{,}175\right)^2 + 10\left(7-7{,}175\right)^2 + 12\left(8-7{,}175\right)^2 \right.
			\\ & \quad \left. + 4\left(9-7{,}175\right)^2 + 2\left(10-7{,}175\right)^2 \right] 
			\\ & \approx 1{,}94.
		\end{align*}
	}
\end{ex}	

\begin{ex}%[0D6H4-4]
	Theo dõi thời gian làm một bài toán (tính bằng phút) của $40$ học sinh, giáo viên lập được bảng sau
	\begin{center}
		\begin{tabular}{|c|c|c|c|c|c|c|c|c|c|c|}
			\hline Thời gian $(x)$	& $4$ & $5$ & $6$ & $7$ & $8$ & $9$ & $10$ & $11$ & $12$ & \\
			\hline Tần số $(n)$		& $6$ & $3$ & $4$ & $2$ & $7$ & $5$ & $5$ & $7$ & $1$ & $N=40$\\
			\hline 
		\end{tabular}
	\end{center}
	Phương sai của mẫu số liệu trên gần với số nào nhất? 
	\choice
	{\True $6$}
	{$12$}
	{$40$}
	{$9$}
	\loigiai{
		Ta có giá trị trung bình của mẫu số liệu là
		\[ \overline{x} = \dfrac{1}{N}\left( n_1x_1+n_2x_2+ \cdots + n_kx_k\right) =\dfrac{317}{40}.\]
		Suy ra phương sai của mẫu số liệu là
		\[ s^2 = \dfrac{1}{N}\left[ n_1\left(x_1-\overline{x}\right)^2 + n_2\left(x_2-\overline{x}\right)^2 + \cdots + n_k\left(x_k-\overline{x}\right)^2\right]  \approx 6.\]
	}
\end{ex}	

\Closesolutionfile{ans}
\indapan{6}{ans/ans-0-B15}

\cauds
\Opensolutionfile{ans}[ans/ans-0-B15-DS]
\begin{ex}%[0D6H4-2]
	Khoảng biến thiên tổng số giờ nắng trong năm của một tỉnh thành được thống kê từ năm 2006 đến 2019 được cho như sau
	\begin{center}
		\begin{tabular}{|c|c|c|c|c|c|c|}
			\hline $1 \, 884$ & $1\, 600$ & $1\, 645$ & $2\, 049{,}9$ & $1\, 913{,}8$ & $1\, 664{,}1$ & $1\, 846{,}5$  \\
			\hline $1\, 964{,}8$ & $1\, 951$ & $2\, 023{,}6$ & $1\, 996{,}2$ & $1\, 699{,}1$ & $1\, 845$ & $2\, 190{,}4$ \\
			\hline 
		\end{tabular}
	\end{center}
	Các mệnh đề sau đúng hay sai?
	\choiceTF
	{Số giờ nắng trung bình trong năm là $1\, 826{,}67$ giờ}
	{\True Số giờ nắng nhỏ nhất $1\, 600$ giờ}
	{\True Số giờ nắng lớn nhất là $2\, 190{,}4$ giờ}
	{Khoảng biến thiên là $520{,}4$}
	\loigiai{
		\begin{itemchoice}
			\itemch Sai. Số giờ nắng trung bình là \[\overline{x}=\dfrac{x_1+x_2+\cdots+x_{14}}{14}=1\, 876{,}67 \text{ giờ}.\]
			\itemch Đúng. Từ bảng số liệu ta có số giờ nắng nhỏ nhất là $1\, 600$ giờ. 
			\itemch Đúng. Từ bảng số liệu ta có số giờ nắng lớn nhất $2\, 190{,}4$ giờ.
			\itemch Sai. Khoảng biến thiên là \[2\, 190{,}4-1\, 600=590{,}4.\]
		\end{itemchoice}
	}
\end{ex}

\begin{ex}%[0D6H4-4]
	Số liệu thống kê tỉ lệ $(\%)$ tốt nghiệp THPT của một địa phương từ năm học 2001 - 2002 đến năm học 2016 - 2017 được cho như sau
	\begin{center}
		\begin{tabular}{|c|c|c|c|c|c|c|c|}
			\hline $98{,}82$ & $97{,}46$ & $99{,}19$ & $98{,}90$ & $98{,}65$ & $79{,}51$ & $85{,}06$ & $86{,}18$  \\
			\hline $98{,}68$ & $99{,}23$ & $99{,}93$ & $99{,}34$ & $99{,}74$ & $93{,}08$ & $97{,}34$ & $97{,}82$ \\
			\hline 
		\end{tabular}
	\end{center}
	Các mệnh đề sau đúng hay sai?
	\choiceTF
	{\True Tỉ lệ tốt nghiệm trung bình $95{,}56 \%$}
	{$99{,}19$ là tỉ lệ $(\%)$ tốt nghiệp THPT cao nhất}
	{Phương sai $s^2\approx 36{,}03$}
	{\True Độ lệch chuẩn $s \approx 6{,}09$}
	\loigiai{
		\begin{itemchoice}
			\itemch Đúng. Tỉ lệ tốt nghiệp trung bình là \[\overline{x}=\dfrac{x_1+x_2+\cdots+x_{16}}{16}=95{,}56 \%.\]
			\itemch Sai. Tỉ lệ tốt nghiệp THPT cao nhất là $99{,}74 \%$.
			\itemch Sai. Phương sai là \[ s^2=\dfrac{1}{16}\left[ \left(x_1-\overline{x}\right)^2+ \left(x_2-\overline{x}\right)^2+\cdots +  \left(x_{16}-\overline{x}\right)^2\right]\approx 37{,}03.\]
			\itemch Đúng. Độ lệch chuẩn là \[s=\sqrt{s^2}\approx 6{,}09.\]
		\end{itemchoice}
	}
\end{ex}

\begin{ex}%[0D6H4-2]
	Cho mẫu số liệu $15$; $20$; $1$; $2$; $4$; $3$; $7$; $5$.
	\\ Các mệnh đề sau đúng hay sai?
	\choiceTF
	{\True Cỡ mẫu là $8$}
	{Khoảng biến thiên của mẫu số liệu là $10$}
	{\True Mẫu số liệu có giá trị trung vị là $4{,}5$}
	{Khoảng tứ phân vị của mẫu là $8$}
	\loigiai{
		Sắp xếp mẫu số liệu theo thứ tự không giảm ta được
		\[1;2;3;4;5;7;15;20.\]
		\begin{itemchoice}
			\itemch Đúng. Từ mẫu số liệu ta có cỡ mẫu là $8$.
			\itemch Sai. Khoảng biến thiên là \[ 20-1=19.\]
			\itemch Đúng. Cỡ mẫu $n=8$ là số chẵn nên trung vị là \[M_e=\dfrac{1}{2}(4+5)=4{,}5.\]
			\itemch Sai. Tứ phân vị thứ nhất và thứ ba của mẫu lần lượt là
			\[Q_1=\dfrac{1}{2}(2+3)=2{,}5 ; Q_3=\dfrac{1}{2}(7+15)=11.\]
			Suy ra khoảng tứ phân vị của mẫu là 
			\[D_Q=Q_3-Q_1=8{,}5.\]
		\end{itemchoice}
	}
\end{ex}

\begin{ex}%[0D6H4-4]
	Điều tra thu nhập của $40$ hộ gia đình ở một bản được số liệu như sau
	\begin{center}
		\begin{tabular}{|c|c|c|c|c|c|c|c|c|c|}
			\hline Thu nhập (triệu đồng/năm) & $4$ & $4{,}5$ & $5$ & $5{,}5$ & $6$ & $6{,}5$ & $7$ & $7{,}5$ & $8$  \\
			\hline Số hộ gia đình & $1$ & $3$ & $4$ & $6$ & $8$ & $7$ & $6$ & $3$ & $2$ \\
			\hline 
		\end{tabular}
	\end{center}
	Các mệnh đề sau đúng hay sai?
	\choiceTF
	{\True Cỡ mẫu là $40$}
	{Mốt của mẫu số liệu là $8$}
	{\True Số trung bình của mẫu số liệu trên là $6{,}1125$}
	{Phương sai của mẫu số liệu trên là $1{,}9561$}
	\loigiai{
		\begin{itemchoice}
			\itemch Đúng. Từ giả thiết ta có cỡ mẫu là $40$.
			\itemch Sai. Có $8$ hộ gia đình thu nhập $6$ triệu đồng/năm nên mốt là $6$.
			\itemch Đúng. Số trung bình của mẫu số liệu trên là \[\overline{x}=\dfrac{4\cdot 1+4{,}5\cdot 3+5\cdot 4+5{,}5\cdot 6+6\cdot 8+6{,}5\cdot 7+7\cdot 6+7{,}5\cdot 3 +8\cdot 2}{40}=6{,}1125.\]
			\itemch Sai. Phương sai của mẫu số liệu trên là 
			\begin{align*}
				s^2 &= \dfrac{1}{40}\left(4^2\cdot 1+4{,}5^2\cdot 3+5^2\cdot 4+5{,}5^2\cdot 6+6^2\cdot 8+6{,}5^2\cdot 7+7^2\cdot 6+7{,}5^2\cdot 3 +8^2\cdot 2\right)-6{,}1125^2
				\\ &\approx 0{,}9561.
			\end{align*}
		\end{itemchoice}
	}
\end{ex}


\Closesolutionfile{ans}
\indapan{3}{ans/ans-0-B15-DS}

\Opensolutionfile{ans}[ans/ans-0-B15-KQ]
\caukq
\begin{ex}%[0D6H4-4]
	Cho mẫu số liệu thống kê $\{1;2;3;4;5;6;7;8;9\}$. Tính độ lệch chuẩn của mẫu số liệu trên (làm tròn đến hàng phần trăm).
	\shortans{$2{,}58$}
	\loigiai{
		Ta có giá trị trung bình của mẫu số liệu là \[\overline{x}=\dfrac{1+2+3+4+5+6+7+8+9}{9}=5.\]
		Suy ra độ lệch chuẩn là \[s=\sqrt{\dfrac{1}{9}\left(1^2+2^2+3^2+4^2+5^2+6^2+7^2+8^2+9^2\right)-5^2}\approx 2{,}58.\]
	}
\end{ex}

\begin{ex}%[0D6H4-4]
	Sản lượng lúa (đơn vị là tạ) của $40$ thửa ruộng thí nghiệm có cùng diện tích được trình bày trong bảng số liệu sau
	\begin{center}
		\begin{tabular}{|c|c|c|c|c|c|c|c|c|c|}
			\hline Sản lượng & $20$ & $21$ & $22$ & $23$ & $24$ &  \\
			\hline Tần số & $5$ & $8$ & $11$ & $10$ & $6$ & $N=40$  \\
			\hline 
		\end{tabular}
	\end{center}
	Tìm phương sai của  mẫu số liệu trên (làm tròn đến hàng phần trăm).
	\shortans{$1{,}54$}
	\loigiai{
		Ta có giá trị trung bình của mẫu số liệu là \[\overline{x}=\dfrac{20\cdot 5+21\cdot 8+22\cdot 11+23\cdot 10+24\cdot 6}{40}=22{,}1.\]
		Suy ra phương sai là \[s^2=\dfrac{1}{40}\left(20^2\cdot 5+21^2\cdot 8+22^2\cdot 11+23^2\cdot 10+24^2\cdot 6\right)-22{,}1^2\approx 1{,}54.\]
	}
\end{ex}

\begin{ex}%[0D6H4-3]
	Cho mẫu số liệu thống kê sau
	\[ -3 \quad 5 \quad 10 \quad 12 \quad 14 \quad 18 \quad 24 \quad 26 \quad 49 \quad 60\]
	Trong các giá trị của mẫu số liệu trên, hãy tìm giá trị ngoại lệ của mẫu số liệu.
	\shortans{$60$}
	\loigiai{
		Từ bảng số liệu ta tìm được số trung vị $Q_2=\dfrac{14+18}{2}=16$, tứ phân vị thứ nhất $Q_1=10$, tứ phân vị thứ ba $Q_3=26$ và khoảng tứ phân vị $\Delta_Q=26-10=16$. Ta có \[\left[ Q_1-1{,}5\Delta_Q;Q_3+1{,}5\Delta_Q\right]=\left[ -14;50\right].\]
		Từ đó ta có $60$ là số liệu ngoại lệ duy nhất của mẫu số liệu.
	}
\end{ex}

\begin{ex}%[0D6H4-4]
	Trong $5$ lần nhảy xa, hai bạn Hùng và Trung có kết quả (đơn vị mét) lần lượt là
	\begin{center}
		\begin{tabular}{|c|c|c|c|c|c|c|c|c|}
			\hline Hùng & $2{,}4$ & $2{,}6$ & $2{,}4$ & $2{,}5$ & $2{,}6$   \\
			\hline Trung & $2{,}4$ & $2{,}5$ & $2{,}5$ & $2{,}5$ & $2{,}6$  \\
			\hline 
		\end{tabular}
	\end{center}
	Hỏi phương sai mẫu số liệu của Hùng gấp của Trung bao nhiêu lần?
	\shortans{$2$}
	\loigiai{
		Phương sai của mẫu số liệu là
		\begin{itemize}
			\item Hùng $s^2=\dfrac{1}{5}\left(2{,}4^2+2{,}6^2+2{,}4^2+2{,}5^2+2{,}6^2\right)-2{,}5^2\approx 0{,}008$.
			\item Trung $s^2=\dfrac{1}{5}\left(2{,}4^2+2{,}5^2+2{,}5^2+2{,}5^2+2{,}6^2\right)-2{,}5^2\approx 0{,}004$.
		\end{itemize}
		Vậy phương sai mẫu số liệu của Hùng gấp của Trung $2$ lần.
	}
\end{ex}

\begin{ex}%[0D6H4-4]
	Bảng số liệu sau thống kê nhiệt độ tại Thành phố Hồ Chí Minh trong một lần đo vào một ngày của năm 2021 như sau
	\begin{center}
		\begin{tabular}{|c|c|c|c|c|c|c|c|c|c|c|c}
			\hline Giờ đo & $1$h & $4$h & $7$h & $10$h & $13$h & $16$h & $19$h & $22$h   \\
			\hline Nhiệt độ ($^\circ \mathrm{C}$) & $27$ & $26$ & $28$ & $32$ & $34$ & $35$ & $30$ & $28$  \\
			\hline 
		\end{tabular}
	\end{center}
	Tìm độ lệch chuẩn của mẫu số liệu đã cho (làm tròn kết quả đến hàng phần trăm).
	\shortans{$3{,}12$}
	\loigiai{
		Ta có số trung bình là \[\overline{x}=\dfrac{27+26+28+32+34+35+30+28}{8}=30^\circ C.\]
		Suy ra độ lệch chuẩn là \[s=\sqrt{\dfrac{1}{8}\left[ (27-30)^2+(26-30)^2+(28-30)^2+\cdots+(28-30)^2\right] }\approx 3{,}12.\]
	}
\end{ex}

\begin{ex}%[0D6H4-4]
	Để biết cây đậu phát triển như thế nào sau khi gieo hạt, bạn Châu gieo $5$ hạt đậu vào $5$ chậu riêng biệt và cung cấp cho chúng lượng nước, ánh sáng như nhau. Sau $2$ tuần, $5$ hạt đậu đã nảy mầm và phát triển thành $5$ cây con. Bạn Châu đo chiều cao từ rễ đến ngọn của mỗi cây (đơn vị mm) và ghi kết quả là mẫu số liệu sau
	\[ 112 \quad 102 \quad 106 \quad 94 \quad 101 \]
	Tìm độ lệch chuẩn của mẫu số liệu trên (kết quả làm tròn đến hàng phần trăm).
	\shortans{$5{,}93$}
	\loigiai{
		Từ bảng số liệu ta có số trung bình là \[ \overline{x}=\dfrac{112+102+106+94+101}{5}=103.\]
		Suy ra ta tính được phương sai và độ lệch chuẩn lần lượt là \[s^2=35{,}2 \text{ và } s\approx 5{,}93.\]
	}
\end{ex}

\Closesolutionfile{ans}
\indapan{6}{ans/ans-0-B15-KQ}